% ============================================================
\chapter{Constrained and Safe RL}
\label{ch:constrained}
% ============================================================

\begin{intuition}
Standard RL maximizes cumulative reward, but real-world applications
require satisfying safety constraints: a robot must avoid collisions,
an autonomous car must obey traffic laws, and an LLM must not generate
harmful content. Constrained and safe RL formalize these requirements
as constraints on cost functions, ensuring that the learned policy is
both performant and safe.
\end{intuition}

% ------------------------------------------------------------
\section{Constrained Markov Decision Processes}
\label{sec:cmdp}
% ------------------------------------------------------------

\begin{definition}[Constrained MDP]
\index{constrained MDP}
A \textbf{Constrained MDP} (CMDP) extends the standard MDP with $K$
cost functions $c_k: \mathcal{S} \times \mathcal{A} \to \R$ and
corresponding thresholds $d_k$:
\[
    \max_\pi J(\pi) = \E_\pi\left[\sum_{t=0}^{\infty} \gamma^t R_t\right]
    \quad \text{s.t.} \quad
    J_{c_k}(\pi) = \E_\pi\left[\sum_{t=0}^{\infty} \gamma^t c_k(S_t, A_t)\right]
    \leq d_k, \quad k = 1, \ldots, K.
\]
\end{definition}

\begin{theorem}[CMDP Optimality]
Under standard regularity conditions, the optimal policy of a CMDP is
a mixture of at most $K+1$ deterministic policies. If the feasible set
is non-empty and the constraint functions are continuous, an optimal
policy exists.
\end{theorem}

\begin{example}[Robot Navigation]
A mobile robot must reach a goal (reward) while limiting the probability
of collision (cost). The CMDP formulation:
\begin{itemize}
    \item Reward: $+1$ for reaching the goal, $-0.01$ per step
    \item Cost: $c(s,a) = \ind_{\{s \in \mathcal{S}_{\text{unsafe}}\}}$
    \item Constraint: $J_c(\pi) \leq 0.05$ (collision rate below 5\%)
\end{itemize}
\end{example}

% ------------------------------------------------------------
\section{Lagrangian Methods}
\label{sec:lagrangian}
% ------------------------------------------------------------

\begin{definition}[Lagrangian Relaxation]
\index{Lagrangian method}
The CMDP is reformulated as an unconstrained saddle-point problem using
Lagrange multipliers $\lambda_k \geq 0$:
\[
    \min_{\lambda \geq 0} \max_\pi \;
    J(\pi) - \sum_{k=1}^{K} \lambda_k \left(J_{c_k}(\pi) - d_k\right)
    = \min_{\lambda \geq 0} \max_\pi \; L(\pi, \lambda).
\]
\end{definition}

\begin{keyformulas}[title=\textbf{Lagrangian Actor-Critic Update}]
\textbf{Policy (primal)}:
\[
    \theta \leftarrow \theta + \alpha_\theta \nabla_\theta \left[
    J(\pi_\theta) - \sum_k \lambda_k J_{c_k}(\pi_\theta)\right]
\]
\textbf{Multipliers (dual)}:
\[
    \lambda_k \leftarrow \max\left(0, \;
    \lambda_k + \alpha_\lambda (J_{c_k}(\pi_\theta) - d_k)\right)
\]
\end{keyformulas}

\begin{algorithme}[title=\textbf{Lagrangian PPO for CMDPs}]
\begin{enumerate}
    \item Initialize policy $\theta$, value networks $V_\phi$, $V_{\phi_k}^{c}$ for each cost,
          multipliers $\lambda_k = 0$
    \item For each iteration:
    \begin{enumerate}
        \item Collect trajectories using $\pi_\theta$
        \item Estimate reward advantages $\hat{A}_t^r$ and cost advantages
              $\hat{A}_t^{c_k}$ using GAE
        \item Update $\theta$ using PPO with modified advantage:
              $\hat{A}_t = \hat{A}_t^r - \sum_k \lambda_k \hat{A}_t^{c_k}$
        \item Update multipliers:
              $\lambda_k \leftarrow \max(0, \lambda_k + \alpha_\lambda(\hat{J}_{c_k} - d_k))$
    \end{enumerate}
\end{enumerate}
\end{algorithme}

\begin{theorem}[Lagrangian Duality for CMDPs]
Under Slater's condition (there exists a strictly feasible policy), strong
duality holds for CMDPs:
\[
    \max_{\pi \in \Pi_C} J(\pi) = \min_{\lambda \geq 0} \max_\pi L(\pi, \lambda)
\]
where $\Pi_C$ is the set of feasible policies. The Lagrangian approach
therefore converges to the optimal constrained policy.
\end{theorem}

% ------------------------------------------------------------
\section{Safe Exploration}
\label{sec:safe_exploration}
% ------------------------------------------------------------

\begin{definition}[Safe Exploration]
\index{safe exploration}
\textbf{Safe exploration} requires that the agent satisfies safety
constraints \emph{during training}, not just at convergence. This is
critical for physical systems where unsafe exploration causes
irreversible damage.
\end{definition}

\begin{definition}[Safety Layer]
A \textbf{safety layer} projects the agent's action onto the safe set
at each step:
\[
    a_{\text{safe}} = \argmin_{a \in \mathcal{A}_{\text{safe}}(s)}
    \norm{a - a_{\text{proposed}}}^2
\]
where $\mathcal{A}_{\text{safe}}(s) = \{a : c(s,a) \leq 0\}$ is the
set of safe actions in state $s$.
\end{definition}

\begin{proposition}[Lyapunov-Based Safety]
A Lyapunov function $L(s) \geq 0$ with $L(s) = 0$ only for safe states
can certify safety. If the policy ensures:
\[
    \E[L(S_{t+1}) | S_t = s, A_t = a] \leq L(s) - \alpha L(s) + \epsilon
\]
for some $\alpha > 0$ and small $\epsilon > 0$, then the system remains
within a safe region with high probability.
\end{proposition}

% ------------------------------------------------------------
\section{Reward Shaping}
\label{sec:reward_shaping}
% ------------------------------------------------------------

\begin{definition}[Potential-Based Reward Shaping]
\index{reward shaping}
Given a potential function $\Phi: \mathcal{S} \to \R$, the shaped reward is:
\[
    R'(s, a, s') = R(s, a, s') + \gamma \Phi(s') - \Phi(s).
\]
\end{definition}

\begin{theorem}[Ng et al., 1999]
Potential-based reward shaping preserves the optimal policy: the set of
optimal policies under $R'$ is identical to that under $R$. Any non-potential
shaping function can change the optimal policy.
\end{theorem}

\begin{remark}
Reward shaping is widely used to incorporate domain knowledge and
accelerate learning. For safety, negative shaping near dangerous states
(e.g., $\Phi(s) = -d(s, \mathcal{S}_{\text{unsafe}})^{-1}$) guides the
agent away from unsafe regions without altering optimality.
\end{remark}

% ------------------------------------------------------------
\section{RLHF: Reinforcement Learning from Human Feedback}
\label{sec:rlhf}
% ------------------------------------------------------------

\begin{definition}[RLHF]
\index{RLHF}
\textbf{Reinforcement Learning from Human Feedback} (Christiano et al.,
2017; Ouyang et al., 2022) trains a policy to align with human
preferences in three stages:
\begin{enumerate}
    \item \textbf{Supervised fine-tuning} (SFT): fine-tune a pretrained
    LLM on demonstration data.
    \item \textbf{Reward modeling}: train a reward model $r_\psi(x, y)$
    from human preference comparisons $(y_w \succ y_l | x)$ using the
    Bradley-Terry model:
    \[
        P(y_w \succ y_l | x) = \sigma\left(r_\psi(x, y_w) - r_\psi(x, y_l)\right)
    \]
    \item \textbf{RL optimization}: optimize the policy $\pi_\theta$ using
    PPO with the learned reward, subject to a KL constraint from the
    reference policy $\pi_{\text{ref}}$:
    \[
        \max_\theta \; \E_{x, y \sim \pi_\theta}\left[r_\psi(x, y)\right]
        - \beta \, D_{\text{KL}}\left(\pi_\theta \| \pi_{\text{ref}}\right)
    \]
\end{enumerate}
\end{definition}

\begin{keyformulas}[title=\textbf{RLHF Objective}]
\[
    J(\theta) = \E_{x \sim \mathcal{D}, y \sim \pi_\theta(\cdot|x)}\left[
    r_\psi(x, y) - \beta \log \frac{\pi_\theta(y|x)}{\pi_{\text{ref}}(y|x)}
    \right]
\]
The KL penalty prevents the policy from diverging too far from the
pretrained model, maintaining coherent language generation.
\end{keyformulas}

\begin{intuition}
RLHF bridges the gap between language model capabilities and human
values. The reward model captures nuanced human preferences that are
difficult to specify programmatically. The KL constraint acts as a
safety mechanism, preventing reward hacking (exploiting quirks of the
learned reward model).
\end{intuition}

\begin{attention}[title=\textbf{Reward Hacking}]
\index{reward hacking}
Without the KL constraint, the policy may find adversarial inputs that
maximize $r_\psi$ without genuinely satisfying human preferences. This
is a form of Goodhart's law: ``when a measure becomes a target, it
ceases to be a good measure.''
\end{attention}

% ------------------------------------------------------------
\section{Direct Preference Optimization (DPO)}
\label{sec:dpo}
% ------------------------------------------------------------

\begin{definition}[DPO]
\index{DPO}
Direct Preference Optimization (Rafailov et al., 2023) eliminates the
explicit reward model by directly optimizing the policy from preferences:
\[
    \mathcal{L}_{\text{DPO}}(\theta) = -\E_{(x, y_w, y_l)}\left[
    \log \sigma\left(\beta \log \frac{\pi_\theta(y_w|x)}{\pi_{\text{ref}}(y_w|x)}
    - \beta \log \frac{\pi_\theta(y_l|x)}{\pi_{\text{ref}}(y_l|x)}\right)
    \right]
\]
\end{definition}

\begin{proposition}[DPO-RLHF Equivalence]
The DPO loss is derived from the closed-form solution of the RLHF
objective under the KL constraint. The optimal policy satisfies:
\[
    \pi^*(y|x) = \frac{\pi_{\text{ref}}(y|x) \exp(r(x,y)/\beta)}{Z(x)}
\]
Substituting this into the Bradley-Terry model yields the DPO loss,
which can be optimized without RL.
\end{proposition}

% ------------------------------------------------------------
\section{Python Example}
\label{sec:constrained:code}
% ------------------------------------------------------------

\begin{codeblock}[title=\textbf{Lagrangian PPO (Simplified)}]
\begin{minted}{python}
import torch

class LagrangianPPO:
    def __init__(self, n_constraints, init_lambda=0.0,
                 lr_lambda=0.01):
        self.lambdas = [init_lambda] * n_constraints
        self.lr = lr_lambda

    def modified_advantage(self, adv_reward, adv_costs):
        """Compute Lagrangian advantage."""
        adv = adv_reward.clone()
        for k, lam in enumerate(self.lambdas):
            adv -= lam * adv_costs[k]
        return adv

    def update_multipliers(self, cost_values, thresholds):
        """Dual ascent on Lagrange multipliers."""
        for k in range(len(self.lambdas)):
            self.lambdas[k] = max(
                0.0,
                self.lambdas[k]
                + self.lr * (cost_values[k] - thresholds[k])
            )
\end{minted}
\end{codeblock}

% ------------------------------------------------------------
\section{Exercises}
\label{sec:constrained:exercises}
% ------------------------------------------------------------

\begin{exercise}[CMDP Formulation]
Formulate the autonomous braking problem as a CMDP. Define the state
space, action space, reward, cost function, and threshold. What makes
this problem challenging compared to standard RL?
\end{exercise}

\begin{exercise}[Lagrangian Method]
Implement Lagrangian PPO on a grid world where the agent must reach a
goal while avoiding certain cells (cost $= 1$ when entering a forbidden
cell). Vary the cost threshold $d$ and observe the trade-off between
reward and constraint satisfaction.
\end{exercise}

\begin{exercise}[Reward Shaping]
Design a potential-based shaping function for the CartPole environment
that encourages keeping the pole upright without changing the optimal
policy. Measure the speedup compared to the unshaped reward.
\end{exercise}

\begin{exercise}[RLHF Pipeline]
Describe the three stages of RLHF for training a dialogue system.
What happens if the reward model is trained on biased preference data?
How can this be mitigated?
\end{exercise}
